\documentclass{article}
\usepackage{amsmath, amssymb, amsthm}
\usepackage{hyperref}
\usepackage{geometry}
\geometry{margin=1in}

\title{Erd\H{o}s Problem \#1109}
\date{Last updated: 2025-12-03}
\author{Source: erdosproblems.com}

\begin{document}
\maketitle

\noindent\textbf{Status:} OPEN \quad \textbf{No prize}

\noindent\textbf{Tags:} number theory

\medskip

\noindent\textbf{Problem Statement:}

Let $f(N)$ be the size of the largest subset $A\subseteq \{1,\ldots,N\}$ such that every $n\in A+A$ is squarefree. Estimate $f(N)$. In particular, is it true that $f(N)\leq N^{o(1)}$, or even $f(N) \leq (\log N)^{O(1)}$?




First studied by Erd\H{o}s and S\'{a}rk\"{o}zy \cite{ErSa87}, who proved\[\log N \ll f(N) \ll N^{3/4}\log N,\]and guessed the lower bound is nearer the truth. S\'{a}rk\"{o}zy \cite{Sa92c} extended this to consider the case of $A+B$ and also looking for sumsets which are $k$-power-free. Gyarmati \cite{Gy01} gave an alternative proof of $f(N)\gg \log N$, and also gave new bounds for the case of $A+B$. Konyagin \cite{Ko04} improved this to\[ \log\log N(\log N)^2\ll f(N) \ll N^{11/15+o(1)}.\]The infinite analogue of this problem is [1103] . (In particular upper bounds for this $f(N)$ directly imply lower bounds for the size of the $a_j$ considered there.)




References

[ErSa87] Erd\H{o}s, P. and S\'{a}rk\"ozy, A., On divisibility properties of integers of the form {$a+a'$} . Acta Math. Hungar. (1987), 117--122.

[Gy01] Gyarmati, Katalin, On divisibility properties of integers of the form {$ab+1$} . Period. Math. Hungar. (2001), 71--79.

[Ko04] Konyagin, S. V., Problems of the set of square-free numbers . Izv. Ross. Akad. Nauk Ser. Mat. (2004), 63--90.

[Sa92c] S\'{a}rk\"ozy, G. N., On a problem of {P}. {E}rd\H{o}s . Acta Math. Hungar. (1992), 271--282.

\noindent\textbf{Related OEIS sequences:} A392164, A392165


\bigskip
\noindent\small{Source: \url{https://www.erdosproblems.com/1109}}
\end{document}
